%%% Бакалаврская ВКР — полный шаблон (все элементы)
%%% Тематика: Машинное обучение, классификация текстов
%%% Компиляция: xelatex → biber → xelatex → xelatex

\documentclass[font=modern]{cmcmsuthesis-bachelor}

\author{Петрова Мария Сергеевна}
\title{Методы машинного обучения для задачи
       классификации текстов}
\subtitle{Выпускная квалификационная работа бакалавра}
\thesissetup{department={Кафедра математических методов прогнозирования}}
\supervisor{name={канд.~физ.-мат.~наук, доцент Козлов~А.Б.}}

\usepackage{biblatex}
\addbibresource{refs.bib}

\begin{document}
\maketitle
\tableofcontents

% ================================================================
\section*{Введение}
\addcontentsline{toc}{section}{Введение}

Задача классификации текстов~\cite{shannon1948} является центральной в NLP.

\textbf{Целью} работы является сравнительный анализ методов классификации.

\textbf{Задачи}:
\begin{enumerate}[beginpenalty=10000]
    \item провести обзор методов;
    \item реализовать модели Байеса, SVM, LSTM;
    \item оценить качество на корпусе текстов;
    \item сформулировать рекомендации.
\end{enumerate}

% ================================================================
\section{Теоретические основы}\label{sec:theory}

\subsection{Основные определения}

Пусть $\vect{x} \in \mathbb{R}^m$~--- вектор признаков документа.

\begin{definition}\label{def:bow}
\emph{Мешком слов} документа $d$ называется вектор
$\vect{x}_d \in \mathbb{R}^m$, где $x_{d,i}$~--- число вхождений
$i$-го слова.
\end{definition}

\begin{definition}\label{def:tfidf}
TF-IDF вес: $\operatorname{tf\text{-}idf}(w,d) =
\operatorname{tf}(w,d)\cdot\log\frac{|D|}{|\{d'\in D\mid w\in d'\}|}$.
\end{definition}

\begin{theorem}[Байесовский классификатор]\label{thm:bayes}
Оптимальный классификатор: $\hat{y} = \arg\max_{c\in\mathcal{Y}}
P(\vect{x}\mid c)\cdot P(c)$.
\end{theorem}

\begin{proof}
Следует из теоремы Байеса и минимизации
$P(\text{ошибка}) = 1 - P(\hat{y}\mid\vect{x})$.
\end{proof}

\begin{corollary}\label{cor:nb}
При условной независимости признаков:
$P(\vect{x}\mid c) = \prod_j P(x_j\mid c)$.
\end{corollary}

\begin{lemma}\label{lem:margin}
Если данные линейно разделимы, то $\exists\,\vect{w}$:
$y_i(\vect{w}^\top\vect{x}_i+b)\geq 1$.
\end{lemma}

\begin{proposition}\label{prop:svmcomplexity}
Обучение SVM решается за $\compl{O}(n^2 m)$.
\end{proposition}

\begin{remark}\label{rem:overfit}
При малой выборке возможно переобучение.
\end{remark}

\begin{example}\label{ex:spam}
Задача фильтрации спама: $\mathcal{Y}=\{\text{спам},\text{не спам}\}$.
\end{example}

\begin{problem}\label{prob:multiclass}
Обобщить бинарный SVM на $|\mathcal{Y}|>2$ классов.
\end{problem}

Команды: $\matr{W}\in\mathbb{F}^{k\times n}_q$,
$\str{S}_n$, $\elt{\sigma}$, $\code{C}$,
$\algo{NaiveBayes}$, $\compl{O}(nm)$, $\crypto{E}$.

% ================================================================
\section{Формулы}\label{sec:formulas}

\subsection{Ненумерованные формулы}
Inline: $x\approx\sin x$ при $x\to 0$.
Display: \[(x_1+x_2)^2=x_1^2+2x_1x_2+x_2^2.\]
Непрерывная дробь:
\[\frac{1}{\sqrt{2}+\displaystyle\frac{1}{\sqrt{2}+\displaystyle\frac{1}{\sqrt{2}+\cdots}}}\]

\subsection{Нумерованные формулы}
\begin{equation}\label{eq:euler}
e=\lim_{n\to\infty}\left(1+\frac{1}{n}\right)^{\!n}.
\end{equation}
\begin{equation}\label{eq:basel}
\lim_{n\to\infty}\sum_{k=1}^n\frac{1}{k^2}=\frac{\pi^2}{6}.
\end{equation}
Ссылки: \eqref{eq:euler} и~\eqref{eq:basel}.

\begin{equation}\label{eq:long}
\begin{multlined}
1+2+3+\dots+50+\dots+\\+96+97+98+99+100=5050.
\end{multlined}
\end{equation}

\subsection{Многострочные формулы}
\[\begin{aligned}
f_W&=\min\!\left(1,\max\!\left(0,\frac{W/W_{\max}}{W_{\text{crit}}}\right)\right),\\
f_T&=\min\!\left(1,\max\!\left(0,\frac{T_s/T_{\text{melt}}}{T_{\text{crit}}}\right)\right).
\end{aligned}\]

\[|x|=\left\{\begin{alignedat}{2}
& &x,\quad&\text{если }x\geqslant 0;\\
&-&x,\quad&\text{если }x<0.
\end{alignedat}\right.\]

\[|x|=\begin{cases}\phantom{-}x,&\text{если }x\geqslant 0;\\-x,&\text{если }x<0.\end{cases}\]

\begin{equation}\label{eq:softmax}
\begin{alignedat}{2}
P(y=c\mid\vect{x})&=\frac{\exp(\vect{w}_c^\top\vect{x})}{\sum_{j=1}^C\exp(\vect{w}_j^\top\vect{x})},&\quad&c=1,\ldots,C.
\end{alignedat}
\end{equation}

\subsection{Субформулы}
\begin{subequations}\label{eq:sub1}
\begin{gather}y=x^2+1\label{eq:s1a}\\y=2x^2-x+1\label{eq:s1b}\\y=3x^2-4x+8\label{eq:s1c}\end{gather}
\end{subequations}
\begin{subequations}\label{eq:sub2}
\begin{align}y&=x^3+x^2+x+1\label{eq:s2a}\\y&=x^2\label{eq:s2b}\end{align}
\end{subequations}
Ссылки: \eqref{eq:s1a}, \eqref{eq:s1b}, \eqref{eq:s2a}, \eqref{eq:s1c}.

Система (Лоренц):
\[\left\{\begin{array}{rl}\dot x=&\sigma(y-x)\\\dot y=&x(r-z)-y\\\dot z=&xy-bz\end{array}\right..\]

Матрица:
\[\left(\begin{array}{ccc}a_{11}&\ldots&a_{1n}\\\vdots&\ddots&\vdots\\a_{n1}&\ldots&a_{nn}\end{array}\right).\]

Длинная inline: $a_1+a_2+a_3+a_4+a_5+a_6+a_7+a_8+a_9+a_{10}+a_{11}+a_{12}+a_{13}+a_{14}+a_{15}+a_{16}+a_{17}+a_{18}+a_{19}+a_{20}=\sum_{i=1}^{20}a_i.$

\subsection{Греческие буквы и алфавиты}
Обычные: $\alpha\beta\gamma\delta\epsilon\theta\lambda\mu\pi\sigma\phi\omega\Gamma\Delta\Omega$.
Прямые: $\symup{\alpha}\symup{\beta}\symup{\gamma}\symup{\delta}\symup{\Gamma}\symup{\Omega}$.
Жирные: $\symbf{\alpha}\symbf{\beta}\symbf{\gamma}\symbf{\Gamma}\symbf{\Omega}$.
Жирные прямые: $\symbfup{\alpha}\symbfup{\beta}\symbfup{\gamma}\symbfup{\Gamma}\symbfup{\Omega}$.

\[\begin{aligned}\mathcal{ABCDEFGHIJKLMNOPQRSTUVWXYZ}\\\mathfrak{ABCDEFGHIJKLMNOPQRSTUVWXYZ}\\\mathbb{ABCDEFGHIJKLMNOPQRSTUVWXYZ}\end{aligned}\]

Русские операторы: $\tan\alpha$/$\cot\beta$/$\csc\gamma$; $\le$/$\leqslant$; $\phi$/$\varphi$; $\epsilon$/$\varepsilon$; $\kappa$/$\varkappa$; $\emptyset$/$\varnothing$.

Отступы:
\[\begin{aligned}\backslash!&\quad f(x)=x^2\!+3x\!+2\\\text{default}&\quad f(x)=x^2+3x+2\\\backslash,&\quad f(x)=x^2\,+3x\,+2\\\backslash;&\quad f(x)=x^2\;+3x\;+2\\\backslash\text{quad}&\quad f(x)=x^2\quad+3x\quad+2\end{aligned}\]

% ================================================================
\section{Форматирование текста}
\begin{description}
\item[Жирный:] \textbf{жирный}.\item[Курсив:] \emph{курсив}.
\item[Жирный курсив:] \textbf{\emph{жирный курсив}}.
\item[Подчёркивание:] \underline{подчёркнутый}.
\item[Зачёркивание:] \sout{зачёркнутый}.
\item[Sans:] \textsf{sans serif font}.
\end{description}

% ================================================================
\section{Таблицы}

\begin{table}[htbp]\caption{Простая таблица}\label{tab:simple}\centering
\begin{tabular}{|l|c|c|c|}\hline Метод&Precision&Recall&F1\\\hline
Байес&0.82&0.79&0.80\\\hline SVM&0.89&0.87&0.88\\\hline
LSTM&0.91&0.90&0.90\\\hline\end{tabular}\end{table}

\begin{table}[htbp]\caption{Booktabs}\label{tab:booktabs}\centering
\begin{tabular}{lrcc}\toprule
\textbf{Метод}&\textbf{Accuracy}&\textbf{Время,с}&\textbf{Модель}\\\midrule
Наивный Байес&0.82&0.3&линейная\\
SVM&0.89&2.1&ядерная\\
LSTM&0.91&45.7&нейросеть\\\bottomrule\end{tabular}\end{table}

\begin{longtable}{@{}rr@{}}\caption{Длинная таблица}\label{tab:long}\\
\toprule Эпоха&Loss\\\midrule\endfirsthead
\toprule Эпоха&Loss\\\midrule\endhead\bottomrule\endfoot
1&2.30\\2&1.85\\3&1.42\\4&1.10\\5&0.87\\
6&0.71\\7&0.58\\8&0.49\\9&0.42\\10&0.37\\\end{longtable}

Ссылка: таблица~\ref{tab:booktabs}.

% ================================================================
\section{Изображения}
%% \begin{figure}[htbp]\centering
%% \includegraphics[width=0.8\textwidth]{images/roc.png}
%% \caption{ROC-кривая}\label{fig:roc}\end{figure}
%%
%% \begin{figure}[htbp]\centering
%% \includegraphics[scale=0.5]{images/confusion.png}
%% \caption{Матрица ошибок (масштаб 0.5)}\label{fig:cm}\end{figure}
%%
%% Ссылка: рисунок~\ref{fig:roc}.

% ================================================================
\section{Списки}
\subsection{Нумерованный}
\begin{enumerate}\item Первый\begin{enumerate}\item Второй\begin{enumerate}\item Третий\begin{enumerate}\item Четвёртый\end{enumerate}\end{enumerate}\end{enumerate}\end{enumerate}
\subsection{Ненумерованный}
\begin{itemize}\item Первый\begin{itemize}\item Второй\begin{itemize}\item Третий\begin{itemize}\item Четвёртый\end{itemize}\end{itemize}\end{itemize}\end{itemize}

% ================================================================
\section{Алгоритмы и листинги}

\begin{algorithm}[H]\caption{Наивный Байес}\label{alg:nb}
\KwInput{выборка $\{(\vect{x}_i,y_i)\}$, тестовый $\vect{x}$}
\KwOutput{метка $\hat{y}$}
\For{каждый класс $c$}{
    Оценить $\hat{P}(c)$\;
    \For{каждый признак $j$}{$\hat{P}(x_j\mid c)$\tcp*{сглаживание}}
}
\Comment{классифицируем}
\uIf{$\max_c\hat{P}(c\mid\vect{x})>\theta$}{\Return{$\hat{y}$}\;}
\uElseIf{два класса равновероятны}{\Return{«неопределённо»}\;}
\Else{\Return{$\arg\max_c\hat{P}(c)\prod_j\hat{P}(x_j\mid c)$}\;}
\end{algorithm}

Ссылка: алгоритм~\ref{alg:nb}.

\begin{lstlisting}[language=Python,caption={Python},label={lst:py}]
from sklearn.naive_bayes import MultinomialNB
from sklearn.metrics import accuracy_score
# Обучение модели
clf = MultinomialNB()
clf.fit(X_train, y_train)
predictions = clf.predict(X_test)
acc = accuracy_score(y_test, predictions)
print(f"Точность: {acc:.3f}")  # Выводим результат
\end{lstlisting}

\begin{lstlisting}[language=C,caption={C},label={lst:c}]
#include <stdio.h>
#include <math.h>
/* Логарифм правдоподобия */
double log_likelihood(int *counts, int n, double alpha) {
    double ll = 0.0;
    // Суммируем логарифмы
    for (int i = 0; i < n; ++i)
        ll += lgamma(counts[i] + alpha);
    return ll;
}
int main(void) { printf("%.4f\n", log_likelihood((int[]){3,5,2}, 3, 1.0)); return 0; }
\end{lstlisting}

\begin{lstlisting}[style=pseudocode,caption={Псевдокод}]
(*@\lstInput@*): выборка $\{(\vect{x}_i,y_i)\}$, тестовый $\vect{x}$
(*@\lstOutput@*): метка класса
procedure Classify($\vect{x}$)
begin
    for each class $c$ do
    begin
        $\log P(c\mid\vect{x}) \gets \log P(c) + \sum_j \log P(x_j\mid c)$
    end
    return $\arg\max_c \log P(c\mid\vect{x})$
end
\end{lstlisting}

\begin{lstpseudocode}[caption={k-NN},label={alg:knn}]
(*@\lstInput@*): выборка, тестовый $\vect{x}$, число соседей $k$
(*@\lstOutput@*): метка класса
function kNN($\vect{x}$, $k$)
begin
    for each $(\vect{x}_i, y_i)$ do
        $d_i \gets \|\vect{x}-\vect{x}_i\|$
    end
    $N_k \gets$ $k$ ближайших
    return голосование по $N_k$
end
\end{lstpseudocode}

% ================================================================
\section{Библиография}
Одиночная~\cite{shannon1948}. Книга~\cite{knuth1984}.
Множественная~\cite{shannon1948,knuth1984,macwilliams1979}.

% ================================================================
\section*{Заключение}
\addcontentsline{toc}{section}{Заключение}
\begin{enumerate}
\item Реализованы три модели классификации.
\item Сравнение в таблице~\ref{tab:booktabs}.
\item LSTM показал лучшее качество.
\end{enumerate}

\clearpage\listoftables\addcontentsline{toc}{section}{Список таблиц}
\printbibliography

\appendix
\section{Исходный код}
\subsection{Модуль предобработки}
Код очистки текстов.
\subsection{Модуль моделей}
Код обучения.

\section{Дополнительные графики}
Дополнительные ROC-кривые.
\subsection{Байес vs SVM}
Данные.

\end{document}
